\section{Underlying Functions and Addresses}
\subsection{Hash Function}

SHRINCS uses SHA-256 truncated to \texttt{n} bytes for all internal hash operations. 
\begin{center}
\texttt{H(m) = SHA-256(m)[0:n]}
\end{center}

\subsection{Tweakable Hash Function}
To prevent multi-target attacks and ensure domain separation between the various components, we employ the tweakable hash function construction defined in FIPS 205 \cite{fips-205}:

\begin{center}
\texttt{Tw(PK.seed, ADRS, m) = H(PK.seed || ADRS || m)}
\end{center}

This usage of tweaked hashing is critical. It ensures that a hash computed at one position in the Merkle tree cannot be substituted for a hash at another position, effectively mitigating generic tree-grafting attacks.

\oleks{\texttt{PK.seed = 16 bytes, ADRS = 32 bytes, m = 16-32 bytes}. SHA-2 operates with 64-byte blocks. So we can add a "0"-padding, like \texttt{Tw(PK.seed, ADRS, m) = H(PK.seed || $\vec{0}$[0,3n] || ADRS || m)} to have the first 64 bytes constant. It will reduce the hashing complexity to 1 block for each tweak function.}

\subsection{Pseudorandom Functions}
We use PRF for deterministic key derivation:

\begin{center}
    \texttt{PRF(PK.seed, ADRS, SK.seed) = Tw(PK.seed, ADRS, SK.seed)}
\end{center}

Additionally we use a pseudorandom function to generate the randomizer \texttt{R} for the randomized hashing of the message to be signed.

\begin{center}
    \texttt{PRF\_msg(SK.prf, PK.seed, opt\_rand, m) = HMAC-SHA-256(SK.prf, PK.seed || opt\_rand || m)},
\end{center}
where \texttt{opt\_rand} is an optional randomness (\texttt{n} bytes), can be \texttt{PK.seed} for deterministic signing.

\subsection{Message Hash Function}
To generate the digest of the message to be signed we use the following function:

\begin{center}
    \texttt{H\_msg(ADRS, R, PK.seed, PK.root, m) = H(ADRS || R || PK.seed || PK.root || m)}
\end{center}

\subsection{Tweak Structure and Domain Separation}
Tweaks are structured as follows to ensure uniqueness across all scheme components:

\begin{center}
\texttt{ADRS = layer || tree\_address || type || words}
\end{center}
where: 
\begin{itemize}
    \item[] \texttt{layer}: (4 bytes) layer in hypertree (\texttt{0} = bottom, \texttt{d-1} = max)
    \item[] \texttt{tree\_address}: (12 bytes) tree address within layer
    \item[] \texttt{type}: (4 bytes) address type
    \item[] \texttt{words}: (12 bytes) type-dependent 12 bytes
\end{itemize}
A total \texttt{ADRS} size is 32 bytes.

\subsubsection{Address Types}
\begin{itemize}
    \item[] \texttt{0x00}: Stateful WOTS+C chains (\texttt{SF\_WOTS\_CHAIN})
    \item[] \texttt{0x01}: Stateful tree nodes (\texttt{SF\_TREE})
    \item[] \texttt{0x02}: Stateless WOTS+C chains (\texttt{SL\_WOTS\_CHAIN})
    \item[] \texttt{0x03}: Stateless tree nodes (\texttt{SL\_TREE})
    \item[] \texttt{0x04}: FORS (\texttt{SL\_FORS})
    \item[] \texttt{0x05}: Message hash (\texttt{H\_MSG})
    \item[] \texttt{0x06}: Pseudorandom function (\texttt{PRF})
    \item[] \texttt{0x07}: Root public key (\texttt{ROOT})
\end{itemize}
Summarizing, \texttt{(0x00, 0x01)} are stateful only types, \texttt{(0x02, 0x03, 0x04)} are stateless only and \texttt{(0x05, 0x06, 0x07)} are shared types.

\subsubsection{Address Manipulation Functions}
\begin{verbatim}
setLayerAddress(ADRS, layer):
    ADRS[0:4] ← toByte(layer, 4)

setTreeAddress(ADRS, tree_addr):
    ADRS[4:16] ← toByte(tree_addr, 12)

setTypeAndClear(ADRS, type):
    ADRS[16:20] ← toByte(type, 4)
    ADRS[20:32] ← toByte(0, 12)

setKeyPairAddress(ADRS, keypair):
    ADRS[20:24] ← toByte(keypair, 4)

setChainAddress(ADRS, chain):
    ADRS[24:28] ← toByte(chain, 4)

setHashAddress(ADRS, hash):
    ADRS[28:32] ← toByte(hash, 4)

setTreeHeight(ADRS, height):
    ADRS[24:28] ← toByte(height, 4)

setTreeIndex(ADRS, index):
    ADRS[28:32] ← toByte(index, 4)
\end{verbatim}

Note that \texttt{setTreeHeight} and \texttt{setChainAddress} both write to bytes \texttt{[24:28]}. These operations are context-dependent so there mustn't be a conflict.

\subsubsection{Type-Dependent Words}

\begin{table}[ht]
\centering
\begin{tabular}{c c c c c}
\hline
\textbf{Type} & \textbf{Meaning} & \textbf{word1 [20:24]} & \textbf{word2 [24:28]} & \textbf{word3 [28:32]} \\
\hline
0x00 & \texttt{SF\_WOTS\_CHAIN} & keypair & chain & hash \\
0x01 & \texttt{SF\_TREE}        & keypair & height & index \\
0x02 & \texttt{SL\_WOTS\_CHAIN} & keypair & chain & hash \\
0x03 & \texttt{SL\_TREE}        & keypair & height & index \\
0x04 & \texttt{SL\_FORS}        & tree\_idx & height & index \\
0x05 & \texttt{H\_MSG}          & 0 & 0 & 0 \\
0x06 & \texttt{PRF}             & W: keypair / F: tree\_idx & W: chain / F: 0 & W: 0 / F: leaf\_idx \\
0x07 & \texttt{ROOT}            & 0 & 0 & 0 \\
\hline
\end{tabular}
\caption{Type-dependent interpretation of ADRS words (bytes [20:32]). W is a WOTS an F is a FORS}
\end{table}

